% Generated by scripts/make_paper_tables.py from the frozen results.
% Do not edit by hand; the next build overwrites it.
\begin{table}[tbp]
\centering
\caption{Model value against the disclosed estimated value, per \$1,000 of principal}
\label{tab:results}
\small
\begin{tabular}{l r r r r r}
\toprule
Note & Model value & s.e. & Disclosed & Gap & Gap (\%) \\
\midrule
N1 & 910.98 & 0.34 & 938.90 & $-$27.92 & $-$2.97 \\
N2\textsuperscript{a} & 894.38 & 0.37 & 943.10 & $-$48.72 & $-$5.17 \\
N3 & 962.56 & 0.29 & 983.40 & $-$20.84 & $-$2.12 \\
N4 & 979.84 & 0.24 & 999.00 & $-$19.16 & $-$1.92 \\
N5 & 960.69 & 0.26 & 972.10 & $-$11.41 & $-$1.17 \\
N6 & 954.01 & 0.30 & 967.30 & $-$13.29 & $-$1.37 \\
\bottomrule
\end{tabular}
\notetext{\textit{Notes:} Model value is the mean discounted payoff over 200,000 paths with a European control variate, at a zero funding spread over the Treasury curve, with the issuer's call exercised by the Longstaff--Schwartz rule. The standard error is of that mean. Gap is model value less disclosed value, so a negative gap means the independent valuation lands below the figure the issuer discloses. Every note shows one. \textsuperscript{a}\,Excluded from the ranges quoted in the text; see Table~\ref{tab:terms}.}
\end{table}
